\scanpage{96}


Pour le vérifier, on note que, pour $x,y$ variables, les deux mbs de la formule sont des formes bil. alternées, et $\{(x,y)\unsure{-|\mapsto}g_x(y)=\alpha^{-1}(x\wedge y)\}$ forme une base du \unsure{Module|module} \unsure{(d'inv.)|(inv.)} des telles formes alternées. Donc on doit avoir

\[ \psi(x)y-\psi(y)x=\unsure{\lambda'|\lambda}\,\alpha^{-1}(x\wedge y)\qquad\text{i.e.} \]

\deleted{avec}

\[ \alpha(\psi(x)y-\psi(y)x)=\lambda\,x\wedge y, \]

\deleted{\unsure{pour}} \unsure{où|\,} $\lambda\in\Gamma O$ est une section \unsure{indépend.|indépend.te} de $x,y$. On veut prouver $\lambda=1$, il suffit donc de le vérifier loc, pour des $x,y$ tels que $x\wedge y$ soit une base de $\det E$. Soit $u$ base de $L$, $x_0$ tel que $\unsure{\varphi|\psi|\Psi}(x_0)=1$, on aura (en prenant $x=u$, $y=x_0$)

\[ \alpha(-\unsure{u|x})=\lambda\underbrace{u\wedge x_0}_{-\alpha(u)} \]

\[ \text{i.e.}\qquad u=\lambda u \]

donc \deleted{(manque de pot !)} $\lambda=1$.

Notons que

\diagram{\begin{tikzcd}[ampersand replacement=\&]0\arrow[r,no head]\&\unsure{O|\mathcal A}\arrow[r,no head]\&\mathcal A\arrow[r,"\varphi",no head]\&L^\vee\arrow[r,no head]\&0\end{tikzcd}}

(transposée de $0-L-E-\underline O-0$)

donne par tensorisation

\diagram{\begin{tikzcd}[ampersand replacement=\&]0\arrow[r,no head]\&L\arrow[r,no head]\&\mathcal A\otimes L\arrow[r,"\psi",no head]\&\underline O\arrow[r,no head]\&0\end{tikzcd}}

et je dis que $g:E-\mathcal A\otimes L$ est un hom d'extensions
% archive sheet number: 95
